\chapter{Configurations des Composants}
\label{annex:configs}

\section{Configuration complète de Mosquitto}

\begin{lstlisting}[language=bash, caption=mosquitto.conf — Configuration TLS/mTLS complète]
# Eclipse Mosquitto 2.0 — Configuration PFE IoT Security
# Port MQTT sécurisé uniquement
listener 8883
protocol mqtt

# TLS/mTLS Configuration
cafile /mosquitto/certs/ca-chain.crt
certfile /mosquitto/certs/server.crt
keyfile /mosquitto/certs/server.key
tls_version tlsv1.3
require_certificate true
use_identity_as_username true

# Sécurité
allow_anonymous false
password_file /mosquitto/config/passwd
acl_file /mosquitto/config/acl.conf

# Performance
max_connections 200
max_inflight_messages 100
max_queued_messages 1000
keepalive_interval 60

# Logging
log_dest file /mosquitto/log/mosquitto.log
log_type error warning notice information
connection_messages true
\end{lstlisting}

\section{Configuration Suricata (extraits)}

\begin{lstlisting}[language=yaml, caption=suricata.yaml — Sections clés]
# Configuration réseau
vars:
  address-groups:
    HOME_NET: "[192.168.10.0/24,192.168.20.0/24]"
    EXTERNAL_NET: "!$HOME_NET"
  port-groups:
    MQTT_PORTS: "8883"

# Mode de capture
af-packet:
  - interface: eth0
    cluster-id: 99
    cluster-type: cluster_flow
    defrag: yes

# Sortie EVE JSON
outputs:
  - eve-log:
      enabled: yes
      filetype: regular
      filename: eve.json
      types:
        - alert
        - http
        - dns
        - tls

# Parser MQTT
app-layer:
  protocols:
    mqtt:
      enabled: yes
      max-msg-length: 1mb
\end{lstlisting}

\section{Configuration MikroTik (Firewall)}

\lstinputlisting[language=bash, caption=mikrotik-config.rsc — Règles de pare-feu]{../../infrastructure/vm/mikrotik-config.rsc}

\section{Configuration Vault (policies)}

\begin{lstlisting}[language=hcl, caption=pki-policy.hcl — Politique Vault pour l'émission de certificats]
# Politique pour l'émission de certificats IoT
path "pki_int/issue/iot-devices" {
  capabilities = ["create", "update"]
}

path "pki_int/issue/services" {
  capabilities = ["create", "update"]
}

path "pki_int/cert/ca" {
  capabilities = ["read"]
}

path "pki/cert/ca" {
  capabilities = ["read"]
}

# Politique de consultation de la CRL
path "pki_int/crl" {
  capabilities = ["read"]
}
\end{lstlisting}

\section{Variables d'environnement requises}

\begin{table}[H]
\centering
\caption{Variables d'environnement du projet}
\label{tab:env_vars}
\begin{tabular}{lll}
\toprule
\textbf{Variable} & \textbf{Description} & \textbf{Exemple} \\
\midrule
\texttt{VAULT\_ADDR} & URL du serveur Vault & \texttt{http://192.168.20.20:8200} \\
\texttt{VAULT\_TOKEN} & Token d'authentification Vault & \texttt{hvs.XXXXXX} \\
\texttt{MQTT\_BROKER} & Adresse du broker Mosquitto & \texttt{192.168.20.10} \\
\texttt{MQTT\_PORT} & Port TLS Mosquitto & \texttt{8883} \\
\texttt{CA\_CERT} & Chemin vers le certificat CA & \texttt{pki/certs/ca/root-ca.crt} \\
\texttt{WAZUH\_MANAGER} & Adresse du manager Wazuh & \texttt{192.168.30.128} \\
\bottomrule
\end{tabular}
\end{table}
